\chapter*{RÉSUMÉ}
\addcontentsline{toc}{chapter}{RÉSUMÉ}
La salinité des sols, très présente au Sine-Saloum, associée à des pratiques d'irrigation et de fertilisation empiriques, compromet la production maraîchère. Ce mémoire, intitulé \textbf{« Optimisation de l’irrigation et de la fertilisation en sols salins par l’agriculture de précision : étude comparative avec les pratiques traditionnelles en cultures maraîchères à Kaolack »}, évalue l'intérêt d'une approche de précision dans ce contexte.

L'expérimentation, effectuée sur le champ d'application Agro TIC de l'USSEIN, compare deux zones de trois plots de $2\,\mathrm{m^2}$ (navet, gombo et laitue). La zone test suit une approche intégrée : étude du sol, amendement ( biochar, PCA, compost et EGB), puis irrigation et fertilisation pilotées à l'aide d'un capteur 8-en-1, de données météorologiques et de la méthode FAO-56 Penman-Monteith \cite{allen_crop_1998-1}. La zone témoin suit la pratique traditionnelle, avec un arrosage fixe. Une application web progressive, AgriSens, conçue avec la méthode MERISE et l'architecture FIWARE, automatise les calculs et la visualisation des données.

Les résultats montrent sur la zone test, un sol plus meuble et plus perméable, sans efflorescence saline, une meilleure croissance (gombo : 41,3~cm contre 21~cm) et une production effective (navet : 4,8~kg sur $2\,\mathrm{m^2}$), alors que dans la zone témoin la production est limitée voire inexistante. Les limites tiennent à l' estimation du NPK à partir de l'EC, à l'absence d'identification des ions, à la faible efficacité des traitements biologiques et à la détection des maladies par simple observation à l'œil nu.

\vspace{1em}
\noindent \textbf{Mots-clés :} agriculture de précision, sols salins, irrigation, fertilisation, évapotranspiration, amendement organique, capteur 8-en-1, cultures maraîchères, AgriSens.